% figures/a1-column-rule.svg — 열 규칙, 그리고 모멘트 항이 곧 팔이라는 것.
%
%   bash _tools/tikz2svg.sh figures/src/a1-column-rule.tex figures/a1-column-rule.svg
%
% 기하는 _tools/a1_column_geo.awk 가 검증한 것과 같다:
%   r = (1.2, 0.8) m, e 는 x_b 에서 -40 deg
%   x e_y - y e_x = -1.3842,  원점에서 작용선까지의 수직거리 = 1.3842 m
% 둘이 같다는 것이 이 그림의 요점이다.
%
% 화면 규약: x_b(선수) 위, y_b(우현) 오른쪽.  TikZ (X,Y) = (y_b, x_b).
\documentclass[border=5pt]{standalone}
% gnc-style.tex — 이 볼트의 TikZ 그림이 공유하는 스타일.
%
% 저널 그림 규약을 스타일 하나로 굳혀 둔다. 그림마다 선 굵기나 화살촉을
% 다시 정하지 않는다 — 그것이 그림이 제각각으로 보이는 원인이다.
%
% 쓰는 법:  % gnc-style.tex — 이 볼트의 TikZ 그림이 공유하는 스타일.
%
% 저널 그림 규약을 스타일 하나로 굳혀 둔다. 그림마다 선 굵기나 화살촉을
% 다시 정하지 않는다 — 그것이 그림이 제각각으로 보이는 원인이다.
%
% 쓰는 법:  % gnc-style.tex — 이 볼트의 TikZ 그림이 공유하는 스타일.
%
% 저널 그림 규약을 스타일 하나로 굳혀 둔다. 그림마다 선 굵기나 화살촉을
% 다시 정하지 않는다 — 그것이 그림이 제각각으로 보이는 원인이다.
%
% 쓰는 법:  \input{gnc-style}  (figures/src/ 안에서)

\usepackage{tikz}
\usepackage{amsmath}
\usepackage{xcolor}
\usetikzlibrary{arrows.meta, positioning, calc, fit, backgrounds, decorations.pathmorphing}

% ---- 색 --------------------------------------------------------------------
% 색은 강조 하나에만 쓴다. 나머지는 검정.
\definecolor{ink}{HTML}{1A1A1A}
\definecolor{accent}{HTML}{7C3AED}
\definecolor{warn}{HTML}{B91C1C}
\definecolor{soft}{HTML}{666666}

% ---- 선과 화살촉 -----------------------------------------------------------
\tikzset{
  % 신호선. 화살촉은 작고 뾰족한 것 하나뿐이다.
  sig/.style   = {ink, line width=0.5pt, -{Stealth[length=4pt, width=3pt]}},
  wire/.style  = {ink, line width=0.5pt},
  % 강조 경로 — 파선으로 구분한다. 흑백 인쇄에서도 살아야 하기 때문이다.
  asig/.style  = {accent, line width=0.5pt, densely dashed,
                  -{Stealth[length=4pt, width=3pt]}},
  awire/.style = {accent, line width=0.5pt, densely dashed},
  % 블록. 흰 바탕, 직각, 검은 테두리.
  blk/.style   = {draw=ink, line width=0.55pt, fill=white,
                  minimum width=17mm, minimum height=8mm, inner sep=2pt, font=\small},
  % 합산점
  sum/.style   = {draw=ink, line width=0.55pt, fill=white, circle,
                  inner sep=0pt, minimum size=4.6mm, font=\scriptsize},
  asum/.style  = {draw=accent, line width=0.55pt, fill=white, circle,
                  inner sep=0pt, minimum size=4.6mm, font=\scriptsize},
  % 분기점
  dot/.style   = {ink, fill=ink, circle, inner sep=0pt, minimum size=1.5mm},
  % 신호 이름 — 선 위에 작게
  sname/.style = {font=\small, inner sep=1.5pt},
  % 그림 안의 짧은 설명. 그림 안의 글은 최소로 둔다
  note/.style  = {font=\scriptsize, soft, inner sep=1.5pt},
  anote/.style = {font=\scriptsize, accent, inner sep=1.5pt},
  % 축
  axis/.style  = {ink, line width=0.5pt},
  tick/.style  = {font=\scriptsize, soft},
}

% 캡션 한 줄 — 그림 아래 가는 선과 함께
\newcommand{\gnccaption}[2]{%
  \node[anchor=north west, text width=#1, align=left, font=\scriptsize, ink]
        at (current bounding box.south west) {\vspace{2pt}#2};}
  (figures/src/ 안에서)

\usepackage{tikz}
\usepackage{amsmath}
\usepackage{xcolor}
\usetikzlibrary{arrows.meta, positioning, calc, fit, backgrounds, decorations.pathmorphing}

% ---- 색 --------------------------------------------------------------------
% 색은 강조 하나에만 쓴다. 나머지는 검정.
\definecolor{ink}{HTML}{1A1A1A}
\definecolor{accent}{HTML}{7C3AED}
\definecolor{warn}{HTML}{B91C1C}
\definecolor{soft}{HTML}{666666}

% ---- 선과 화살촉 -----------------------------------------------------------
\tikzset{
  % 신호선. 화살촉은 작고 뾰족한 것 하나뿐이다.
  sig/.style   = {ink, line width=0.5pt, -{Stealth[length=4pt, width=3pt]}},
  wire/.style  = {ink, line width=0.5pt},
  % 강조 경로 — 파선으로 구분한다. 흑백 인쇄에서도 살아야 하기 때문이다.
  asig/.style  = {accent, line width=0.5pt, densely dashed,
                  -{Stealth[length=4pt, width=3pt]}},
  awire/.style = {accent, line width=0.5pt, densely dashed},
  % 블록. 흰 바탕, 직각, 검은 테두리.
  blk/.style   = {draw=ink, line width=0.55pt, fill=white,
                  minimum width=17mm, minimum height=8mm, inner sep=2pt, font=\small},
  % 합산점
  sum/.style   = {draw=ink, line width=0.55pt, fill=white, circle,
                  inner sep=0pt, minimum size=4.6mm, font=\scriptsize},
  asum/.style  = {draw=accent, line width=0.55pt, fill=white, circle,
                  inner sep=0pt, minimum size=4.6mm, font=\scriptsize},
  % 분기점
  dot/.style   = {ink, fill=ink, circle, inner sep=0pt, minimum size=1.5mm},
  % 신호 이름 — 선 위에 작게
  sname/.style = {font=\small, inner sep=1.5pt},
  % 그림 안의 짧은 설명. 그림 안의 글은 최소로 둔다
  note/.style  = {font=\scriptsize, soft, inner sep=1.5pt},
  anote/.style = {font=\scriptsize, accent, inner sep=1.5pt},
  % 축
  axis/.style  = {ink, line width=0.5pt},
  tick/.style  = {font=\scriptsize, soft},
}

% 캡션 한 줄 — 그림 아래 가는 선과 함께
\newcommand{\gnccaption}[2]{%
  \node[anchor=north west, text width=#1, align=left, font=\scriptsize, ink]
        at (current bounding box.south west) {\vspace{2pt}#2};}
  (figures/src/ 안에서)

\usepackage{tikz}
\usepackage{amsmath}
\usepackage{xcolor}
\usetikzlibrary{arrows.meta, positioning, calc, fit, backgrounds, decorations.pathmorphing}

% ---- 색 --------------------------------------------------------------------
% 색은 강조 하나에만 쓴다. 나머지는 검정.
\definecolor{ink}{HTML}{1A1A1A}
\definecolor{accent}{HTML}{7C3AED}
\definecolor{warn}{HTML}{B91C1C}
\definecolor{soft}{HTML}{666666}

% ---- 선과 화살촉 -----------------------------------------------------------
\tikzset{
  % 신호선. 화살촉은 작고 뾰족한 것 하나뿐이다.
  sig/.style   = {ink, line width=0.5pt, -{Stealth[length=4pt, width=3pt]}},
  wire/.style  = {ink, line width=0.5pt},
  % 강조 경로 — 파선으로 구분한다. 흑백 인쇄에서도 살아야 하기 때문이다.
  asig/.style  = {accent, line width=0.5pt, densely dashed,
                  -{Stealth[length=4pt, width=3pt]}},
  awire/.style = {accent, line width=0.5pt, densely dashed},
  % 블록. 흰 바탕, 직각, 검은 테두리.
  blk/.style   = {draw=ink, line width=0.55pt, fill=white,
                  minimum width=17mm, minimum height=8mm, inner sep=2pt, font=\small},
  % 합산점
  sum/.style   = {draw=ink, line width=0.55pt, fill=white, circle,
                  inner sep=0pt, minimum size=4.6mm, font=\scriptsize},
  asum/.style  = {draw=accent, line width=0.55pt, fill=white, circle,
                  inner sep=0pt, minimum size=4.6mm, font=\scriptsize},
  % 분기점
  dot/.style   = {ink, fill=ink, circle, inner sep=0pt, minimum size=1.5mm},
  % 신호 이름 — 선 위에 작게
  sname/.style = {font=\small, inner sep=1.5pt},
  % 그림 안의 짧은 설명. 그림 안의 글은 최소로 둔다
  note/.style  = {font=\scriptsize, soft, inner sep=1.5pt},
  anote/.style = {font=\scriptsize, accent, inner sep=1.5pt},
  % 축
  axis/.style  = {ink, line width=0.5pt},
  tick/.style  = {font=\scriptsize, soft},
}

% 캡션 한 줄 — 그림 아래 가는 선과 함께
\newcommand{\gnccaption}[2]{%
  \node[anchor=north west, text width=#1, align=left, font=\scriptsize, ink]
        at (current bounding box.south west) {\vspace{2pt}#2};}


\begin{document}
\begin{tikzpicture}[x=1mm, y=1mm]

\def\S{22}                 % 1 m = 22 mm
\def\PXt{17.6}             % 추진기 TikZ x = 0.8 * S
\def\PYt{26.4}             % 추진기 TikZ y = 1.2 * S
\def\FXt{23.33}            % 수선의 발  = 1.0603 * S
\def\FYt{19.58}            %            = 0.8899 * S

% ---- 몸체 좌표계 -----------------------------------------------------------
\draw[axis, -{Stealth[length=4pt,width=3pt]}] (0,0) -- (0,16);
\draw[axis, -{Stealth[length=4pt,width=3pt]}] (0,0) -- (16,0);
\node[font=\small, anchor=east]  at (-1.2,15) {$x_b$};
\node[font=\small, anchor=north] at (15,-1.2) {$y_b$};
\fill[ink] (0,0) circle (0.7);
\node[font=\small, anchor=north east] at (-0.8,-0.8) {$o_b$};

% ---- 작용선 (양쪽으로 연장) ------------------------------------------------
\draw[soft, line width=0.35pt, densely dashed]
      ({\PXt+0.643*17}, {\PYt-0.766*17}) -- ({\PXt-0.643*23}, {\PYt+0.766*23});
\node[note, anchor=west] at ({\PXt+0.643*17+1.5}, {\PYt-0.766*17-2.5}) {the line of action};

% ---- 위치벡터 r ------------------------------------------------------------
\draw[ink, line width=0.8pt, -{Stealth[length=4.5pt,width=3.4pt]}]
      (0,0) -- (\PXt,\PYt);
\node[font=\small, anchor=south east] at ({0.55*\PXt}, {0.55*\PYt}) {$\mathbf{r}$};

% ---- 모멘트 팔 : 원점에서 작용선에 내린 수선 -------------------------------
\draw[accent, line width=0.8pt, -{Stealth[length=4.5pt,width=3.4pt]}]
      (0,0) -- (\FXt,\FYt);
\fill[accent] (\FXt,\FYt) circle (0.55);
% 직각 표시
\draw[accent, line width=0.4pt]
      ($(\FXt,\FYt)+(-1.643,-1.107)$) -- ($(\FXt,\FYt)+(-0.877,-2.873)$)
   -- ($(\FXt,\FYt)+(0.766,-1.766)$);
% 팔 라벨은 작도 아래로 빼고 지시선을 단다. 작용선 라벨과 겹치지 않게.
\draw[accent, line width=0.3pt, densely dotted] (11.7,9.8) -- (10,-3.5);
\node[font=\scriptsize, accent, anchor=north] at (10,-4)
     {the moment arm, $1.3842$\,m};

% ---- 추진기와 추력 ---------------------------------------------------------
\fill[ink] ({\PXt-1.1},{\PYt-1.1}) rectangle ({\PXt+1.1},{\PYt+1.1});
\node[note, anchor=west] at ({\PXt+2}, {\PYt+1}) {the thruster};
\draw[accent!0!ink, line width=0.9pt, -{Stealth[length=5pt,width=3.6pt]}]
      (\PXt,\PYt) -- ({\PXt-0.643*22}, {\PYt+0.766*22});
\node[font=\small, anchor=south east] at ({\PXt-0.643*22-0.5}, {\PYt+0.766*22})
     {$T\mathbf{e}$};

% ---- e 의 성분 -------------------------------------------------------------
\draw[soft, line width=0.35pt, densely dotted]
      (\PXt,\PYt) -- (\PXt, {\PYt+0.766*22});
\draw[soft, line width=0.35pt, densely dotted]
      (\PXt, {\PYt+0.766*22}) -- ({\PXt-0.643*22}, {\PYt+0.766*22});
\node[font=\scriptsize, anchor=west] at ({\PXt+0.8}, {\PYt+0.766*11}) {$e_x$};
\node[font=\scriptsize, anchor=south] at ({\PXt-0.643*11}, {\PYt+0.766*22+0.5}) {$e_y$};

% ===================== 오른쪽 : 규칙 =======================================
\begin{scope}[shift={(62,0)}]
  \node[font=\small\bfseries, anchor=west] at (0, 50) {the rule};

  \node[anchor=north west, align=left, text width=72mm, font=\small] at (0, 46) {%
    one newton from a thruster at $(x,y)$ pointing along $\mathbf{e}$:};
  \node[anchor=north, font=\normalsize] at (36, 39) {%
    $\mathbf{b} = [\,e_x,\ \ e_y,\ \ x\,e_y - y\,e_x\,]^{\top}$};
  \node[anchor=north west, align=left, text width=72mm, font=\scriptsize, soft] at (0, 31) {%
    and that column depends on geometry alone.};

  \node[anchor=north west, align=left, text width=72mm, font=\small] at (0, 25) {%
    $n$ thrusters, $n$ columns:\quad
    $\boldsymbol{\tau} = [\,X\ \ Y\ \ N\,]^{\top} = \mathbf{B}\mathbf{f}$};

  \node[anchor=north west, align=left, text width=72mm, font=\scriptsize] at (0, 16) {%
    \textbf{The check that makes the rule believable.}
    With $x=1.2$\,m, $y=0.8$\,m and $\mathbf{e}$ $40^\circ$ to port of the bow,
    $x e_y - y e_x = \mathbf{-1.3842}$, and the violet arm measures
    $\mathbf{1.3842}$\,m. The third entry is not a new idea — it is the ordinary
    \emph{force times perpendicular distance}, written so the sign comes out on
    its own.\\[4pt]
    \textbf{The base Otter, from the same rule.} Both propellers face forward, so
    $\mathbf{e}=(1,0)$, and they sit at $x=0$, $y=\mp0.395$\,m, giving
    $x e_y - y e_x = \pm 0.395$. Here $e_y = 0$ for both, so the middle row of
    $\mathbf{B}$ is empty — and the Otter can never be pushed sideways.\\[4pt]
    \textcolor{warn}{\textbf{$\mathbf{e}$ must be a unit vector.}} Entering $(2,0)$
    instead of $(1,0)$ silently doubles that thruster's authority, and every rank
    test still passes.};
\end{scope}

\end{tikzpicture}
\end{document}
